\documentclass[12pt]{article}
\usepackage{samstyle} 

\newcommand{\NuevaPagina}[1]{\newpage
\mbox{
\begin{minipage}{3cm}
    \includegraphics[width=3cm]{logo}
\end{minipage}
\begin{minipage}{14.7cm}
  \begin{flushright}
\hfill P�gina #1 \hfill Primer parcial-ST0240 \hfill Marzo
\hfill C�digo: \rule{4cm}{0.3pt}
  \end{flushright}
\end{minipage}}\vspace{3pt}

\rule{18cm}{1.5pt}}

%%% COMIENZO DE DOCUMENTO %%%
\begin{document}
\pagestyle{empty}
%%% Encabezado
\mbox{
    \begin{minipage}{5cm}
        \includegraphics[width=5cm]{logo}
    \end{minipage}
    \begin{minipage}{13cm}
    \begin{center}
        \textsl{
        \textbf{\Large Escuela de Ingenier�as}\\
        \textbf{\large Programaci�n de Computadores - ST0240}\\
        \textbf{\large Pr�ctica 1 (Valor 20\%)}}
    \end{center}
    \end{minipage}
}
\begin{flushright}
  \begin{minipage}{4.5cm}
    Nota: \hrulefill
  \end{minipage}
\end{flushright}\vspace{-2mm}

% Nombre: \rule{9cm}{0.3pt} \; C�digo: \hrulefill\vspace{6pt}

Profesor: \underbar{\hspace{1cm} Sergio Monsalve
\hspace{3.1cm}} \; \hfill Grupo: \underbar{\hspace{1cm} 063 } \;
Marzo \underbar{\hspace{1cm}} de 2016 \rule{18cm}{1.5pt}
%%% Fin del encabezado

La siguiente es la descripci�n de la Practica 1 del curso de Programaci�n de computadores, esta debe ser entregada m�ximo el d�a domingo 3 de Abril a las 11:59pm a trav�s de la plataforma de EAFIT Interactiva.\\

Suponga que usted es el profesor(a) de un curso de Programaci�n de Computadores, tal curso tiene dos parciales del 15\%, dos practicas del 20\% y una pr�ctica final del 30\%. \\

Como entradas usted tiene un archivo <<indece.txt>> que contiene la cantidad de archivos, la cantidad de estudiantes, el nombre de cada archivo, la cantidad de estudiantes que presetaron esa actividad evaluativa, que estudiantes no presentaron tal evaluaci�n, la cantidad de puntos que tenia cada actividad y los porcentajes de cada punto. Un archivo por cada evaluaci�n donde esta el identificador de cada estudiante (entero de 1 a 27) y las notas por cada punto. 

Se requiere que usted elabore un programa que cumpla con los siguientes requisitos: 

\begin{enumerate}
\item \textbf{(Valor 15\%)} Crear una interfaz de linea de comandos que permita: (Entrada y Salida por teclado)
\end{enumerate}

\begin{itemize}
  \item Ingresar la cantidad de archivos a analizar.
  \item Ingresar la cantidad de Estudiantes.
  \item Ingresar el nombre de los archivos. 
\end{itemize}
\vspace{0.5cm}

\begin{enumerate}[resume]
\item \textbf{(Valor 10\%)} Controlar los siguientes casos mediante excepciones:
\end{enumerate}

\begin{itemize}
  \item En caso de que no coincida alguno de los nombres, pedirle al usuario que lo ingrese nuevamente.
  \item Verificar que la suma de los porcentajes de ning�n parcial exceda el porcentaje m�ximo (100).
  \item verificar que no haya ning�n archivo vac�o.
  \item verificar que no se hacen divisiones por 0.
\end{itemize}
\vspace{0.5cm}

\begin{enumerate}[resume]
\item \textbf{(Valor 20\%)} Manipulaci�n datos:
\end{enumerate}

\begin{itemize}
  \item Realizar un arreglo con las notas de cada uno de los estudiantes
  \item Realizar un arreglo por estudiante con las notas de cada evaluaci�n
\end{itemize}
 \vspace{0.5cm}

\begin{enumerate}[resume] 
\item \textbf{(Valor 20\%)}Procesamiento
\end{enumerate}

\begin{itemize}
  \item Encuentre el Promedio de las notas de cada estudiante.
  \item Encuentre la desviaci�n t�pica para cada estudiante.
  \item Encuentre la Nota promedio para cada evaluaci�n
  \item Encuentre la desviaci�n t�pica por evaluaci�n.
\end{itemize}
\vspace{0.5cm}

\begin{enumerate}[resume]
\item \textbf{(Valor 15\%)} Manejo de bibliotecas: Matplotlib
\end{enumerate}


\begin{itemize}
  \item Graficar el Progreso del Estudiante Dado (notas de las evaluaciones)
  \item Graficar el progreso de todos los estudiantes (Gr�fica de lineas con linea por cada estudiante)
\end{itemize}
\vspace{0.5cm}

\begin{enumerate}[resume]
\item \textbf{(Valor 20\%)} Interpretaci�n y Reportes
\end{enumerate}


\begin{itemize}
  \item Determinar cu�les estudiantes han mejorado su rendimiento y mostrarlo mediante una gr�fica (solo estudiantes que siempre mejoraron su nota)
\end{itemize}


\section{Ap�ndice}

\subsection{Indice.txt} % (fold)
\label{sub:indice_txt}

% subsection indice_txt (end)
\lstinputlisting{Indice.txt}

\newpage

\subsection{Parcial1.txt} % (fold)
\label{sub:parcial1_txt}

\lstinputlisting{Solucion/Parcial2.txt}
% subsection parcial1_txt (end)
\end{document}
